\documentclass[11pt]{article}
\usepackage{amsmath,amssymb,amsthm,geometry}\geometry{margin=1in}
\newtheorem{theorem}{Theorem}\newtheorem{lemma}{Lemma}
\title{Closed progressive content spans and the separate causal lift}
\date{17 November 2026}\begin{document}\maketitle
\paragraph{Audit.} Independently typed content closure, not a paid-module theorem. The pure base here is small categories with all functors; Sets embed as discrete categories with all functions, including noninjective ones. Progressive occurrence/atomic signatures remain additional retained decorations, and their compatible arrows form a restricted interface; arbitrary pure functors are not claimed to preserve a paid atomic signature. The construction repairs the changing source domain of pointed content substitution by retaining an apex and its map to the original source. It does not choose a path in that apex. All content collections are entire weighted ordered collections in the original symmetric monoidal content doctrine.

\paragraph{Source route and base distinction.} Companion files in this directory are \texttt{\detokenize{Atomic_Protocol_Associator_and_Pure_Base_Boundary.tex}}, \texttt{\detokenize{Progressive_Graded_Spans_and_Causal_Lift.tex}}, and \texttt{\detokenize{Owner_Transparent_Content_Activation_Face.tex}}. Content spans close over small categories and all functors; this does not assert closure of the separately decorated atomic/pure interface or paid causal actions. The prior constructions are \texttt{\detokenize{Progressive_Content_Substitution_Boundary.tex}}, \texttt{\detokenize{Owned_Random_Response_Request_Assembly.tex}}, and \texttt{\detokenize{Causal_Temporal_and_Simultaneous_Request_Assembly.tex}}.

\section{Progressive polynomial actions on the pure base}
For a small category $B$ define $P_CB$ to have objects $u(b_1,\ldots,b_n)$ and $d(c;b_1,\ldots,b_n)$ with $c\in C_n$. An arrow consists of an order-preserving injection of positions and a $B$ arrow at each retained position. An unfinished object may map to either phase. From a completed object the target must be completed with the same $c$, and the position map is the identity. Thus completed objects may have nonidentity component arrows when $B$ is not discrete. This is a deliberate extension of the old progressive observation category: for discrete $B$ it is exactly $\mathcal D_C(B)$. The phase is parent completion, not an assertion that arbitrary category objects $b_i$ are completed processes.

A functor $f:B\to B'$ acts componentwise. Full collection maps act on every completed shape and retain their actual coefficient arrows independently of these raw categories. Define
\[
 \eta_B:B\to P_I B,\quad b\mapsto d(1;b),\qquad
 \mu_{C,D,B}:P_C P_D B\to P_{C\circ D}B.
\]
The second flattens ordered lists. Its output is completed exactly when the outer object and every inner object are completed, with the ordinary grafted shape; otherwise it is unfinished. Arrows flatten their nested injections and component arrows. This is well typed: a completed flattened source forces all its phases completed, hence none of these levels can add positions or change shape. Otherwise an unfinished source permits the flattened injection.

The multiplication is generally noninjective. If the outer shape is known but some inner object is unfinished, the output forgets that known shape, the grouping and empty inner slots. No paid actor or owner is therefore recoverable from the flattened output alone. These data remain in the span apex below.

\begin{lemma}[Cartesian graded action]
Each $P_C$ preserves strict pullbacks of categories. The natural transformations $\eta$ and $\mu_{C,D}$ are cartesian with respect to pure functors. Multiplication is associative up to the full collection associator and has the full collection unit laws.
\end{lemma}
\begin{proof}
A pair of progressive tuples with equal images has the same phase, shape when present, number of positions and position injection on each arrow. Its components therefore lift uniquely into the componentwise pullback category. This proves preservation of strict pullbacks on objects and arrows.

For cartesianness of $\mu$, fix a functor $f:B\to B'$, a nested tuple in $P_C P_D B'$ and a flattened tuple in $P_{C\circ D}B$ whose images agree. The nested tuple supplies every group boundary and phase, including empty groups. Split the flattened $B$ tuple at precisely those boundaries and retain those phases: this is the unique nested $B$ lift. On arrows, the given nested injection and flattened component arrows split in the same unique way. This proves that the naturality square is a strict pullback; it does not assert that $\mu$ itself is injective. The unit argument is the one-position special case.

For three levels the final phase is completed exactly when every phase in the nested tree is completed. In that case both routes give its full grafted shape, and otherwise both give the same flattened unfinished list. Both routes carry exactly the same component arrows and ordered injection. The full collection associator identifies the completed shapes. The unit equations follow by inserting completed unary singleton tuples. Their coefficient equations are ordinary full substitution and symmetric monoidal coherence at every shape, not just reached shapes.\end{proof}

\section{Content spans and actual composition}
A content arrow from $A$ to $B$ is a triple $(C,K,p,q)$
\[
                  A\xleftarrow{p}K\xrightarrow{q}P_CB,
\]
with $K$ small. Neither leg chooses a stochastic law. Its full grade is $C$. For a second arrow $(D,L,r,s)$ from $B$ to $E$, define the composite apex by the actual strict pullback
\[
 W=K\times_{P_CB}P_C L,
 \qquad q(k)=P_Cr(\ell),
\]
and use the legs
\[
 A\xleftarrow{p\pi_K}W
 \xrightarrow{\mu_{C,D,E}\,P_Cs\,\pi_{P_CL}}
 P_{C\circ D}E.                                      \tag{1}
\]
Thus each current $B$ occurrence is lifted into an actual object of $L$, and each retained $B$ arrow into an actual arrow there. Equality in the pullback retains those lifts, rather than assuming they exist uniquely. Empty fibres are permitted at this content-span level.

Identity is $A\xleftarrow{1}A\xrightarrow{\eta_A}P_IA$. A vertical comparison has an apex functor over the source, a full collection map and the corresponding strict right-body square; all actual collection coefficient maps remain fields. Restricting to a fixed grade gives the ordinary maps of these spans. More general lax body comparisons would require a separately specified 2-cell convention and are not asserted here.

\begin{theorem}[Closed content-span composition]
Formula (1) is associative by a canonical apex isomorphism, with the displayed identities as units. These isomorphisms satisfy the pentagon and triangle and retain the entire full-grade maps. Composition is functorial on the stated strict vertical comparisons. Thus it supplies a graded span bicategorical content composition over one pure base, independently of paid contexts.
\end{theorem}
\begin{proof}
For three arrows the two apex expressions identify the same first apex object, one second-apex object at each of its positions, and one third-apex object at each position exposed by those second objects. Their arrows retain the same lifts at every level. The cartesian multiplication lemma gives precisely the unique splitting needed when one route first flattens the observation; the first two apex objects still retain every latent group, including groups with no exposed output. Pullback preservation gives the other rebracketing. Hence the two expressions are canonically isomorphic, with the same original-source projection and the same final body after the full-grade associator. For four arrows both routes keep the identical nested lift tree, proving the pentagon on actual objects and arrows. Cartesian unit squares give the unit apex isomorphisms and triangle. Full coefficient coherence is the full collection coherence, independent of the apex supports. Componentwise apex functors induce pullback maps, and naturality of $P$ and $\mu$ gives functoriality of strict comparisons.\end{proof}

\paragraph{Pure actions and static recovery.}
Contravariance in a pure source functor $a:A'\to A$ is pullback of $p$ along $a$. Covariance in a pure target functor $b:B\to B'$ is postcomposition by $P_Cb$. These are defined for all functors, with the usual pullback comparison isomorphisms. For discrete sets, embed an original static content body $n:A\to F_CB$ by $K=A$, $p=1$ and $q(a)=d(c_a;\mathbf b_a)$. Composing two such embedded arrows yields an apex canonically $A$, because each component lift through the second identity leg is forced. Formula (1) is exactly static content substitution, with its full grade. Strict vertical comparisons between such identity-left discrete arrows recover the original full-grade/body squares. This is an exact embedded static subinterface, not an assertion that every progressive span has such a normal form.

A pointed child protocol construction is represented by a nontrivial apex with source projection to the parent protocol. Its changing dependent domain is now retained by the left leg, rather than mistaken for a fixed-source body. Marks, occurrence functors and primitive atomic scope annotations may be carried on these apexes, but closure of a chosen decorated subclass requires its own compatibility theorem. The reviewed atomic associator applies when both routes retain the same primitive addresses and transported relations; arbitrary span composition is not claimed to preserve those annotations automatically.

\section{Why this closure does not choose a paid execution}
To drive $(p:K\to A)$ from a paid source observation $o:\mathcal H\to A$, where $\mathcal H$ retains full raw histories, a deterministic causal realization needs an actual lift $\widetilde o:\mathcal H\to K$ with $p\widetilde o=o$, plus the required initial/paid/raw/coefficient fields. The span alone supplies only the pullback category of possible lifts. A discrete opfibration $p$, together with a chosen lift of each start, is a sufficient unique-path-lift condition for the raw history functor; it does not itself supply any paid law or coefficient map. General opfibrations with chosen cocartesian lifts supply a narrower chosen transport, but arbitrary vertical motion remains extra data.

Two small examples obstruct automatic realization. If an input arrow has no lift from the selected apex object, a source paid transition cannot be realized. If $A$ is terminal and $K$ is a discrete two-object category, both apex objects lie over the same input and the span chooses neither. Initializing a distribution on them requires additional law data; interpreting that choice as random scheduling would change the deterministic control interface. With infinite fibres even the finite raw-row hypothesis can fail. Full off-policy and zero histories need lifts too, not just positive paths. Noninjective pure maps may merge histories, so complete-control descent additionally requires agreement on their mapped complete commands.

Consequently (1) repairs content-domain closure but not the paid crossing by itself. A paid action must retain the actual causal lift or its admitted family of joint request fillings, the same-event owner/predecessor constraints, old/new actor ages, actual scope coefficient arrows and full poststate/body squares. Temporal concatenation must retain separate paid events and their chronological factors. Forgetting those fields and then taking a category coend cannot reconstruct them. The next mixed-face theorem must establish closure of that explicitly decorated span/action class, or exhibit its failure; the abstract span composition alone is not that theorem.
\end{document}
